\documentclass[a4paper]{article}
\usepackage{amsfonts}
\usepackage{amsmath}
\usepackage{amssymb}
\usepackage{graphicx}     

\newcommand{\Bgtan}{\operatorname{Bgtan}}
\newcommand{\Bgsin}{\operatorname{Bgsin}}

\begin{document}
  
\noindent \large Auteur: Zaky Souai \\
\noindent \large Studierichting:  Informatica\\
\noindent \large Oefeningenreeks 4A nr. 42.15\\

\bigskip

\normalsize

$\displaystyle \int \dfrac{x^3+x}{x^2-1}\,dx$\\

$\displaystyle \dfrac{x^3+x}{x^2-1}=x+\dfrac{2x}{x^2-1}$\\

$\displaystyle x^2-1=(x-1)(x+1)$\\

$\displaystyle \dfrac{2x}{x^2-1}=\dfrac{2x}{(x-1)(x+1)}=\dfrac{A}{x-1}+\dfrac{B}{x+1}$\\

$\displaystyle A=\left.\dfrac{2x}{x+1}\right|_{x=1}=\dfrac{2}{2}=1$\\

$\displaystyle B=\left.\dfrac{2x}{x-1}\right|_{x=-1}=\dfrac{-2}{-2}=1$\\

$\displaystyle \dfrac{x^3+x}{x^2-1}=x+\dfrac{1}{x-1}+\dfrac{1}{x+1}$\\

$\displaystyle \int \dfrac{x^3+x}{x^2-1}\,dx=\int x\,dx+\int\dfrac{dx}{x-1}+\int\dfrac{dx}{x+1}$\\

$\displaystyle =\dfrac{x^2}{2}+\ln|x-1|+\ln|x+1|+c$\\


\begin{thebibliography}{999}
%\bibitem{a} Referentie
\end{thebibliography}
\end{document}